\uuid{FmPI}
\niveau{L1}
\module{Algèbre}
\chapitre{Injection, surjection, bijection}
\sousChapitre{Application}
\cpge{mpsi-mp2i.raisonnement-vocabulaire-ensembliste.c, pcsi.raisonnement-vocabulaire-ensembliste.d, ptsi.raisonnement-vocabulaire-ensembliste.d, bcpst1.fonctions-reelles-usuelles.a, bcpst1.vocabulaire-applications}
%!TeX root=../../../encours.nouveau.tex
%%% Début exercice %%%

\duree{5}
\difficulte{1}
\auteur{Antoine Crouzet}
\datecreate{2024-12-01}
\titre{Des puissances de $-1$}
\contenu{
\question{Soit $f:\Z\dans \R$ définie pour tout entier $x$ par $f(x)=(-1)^x$. Déterminer l'image de la fonction $f$, $f(\Z)$. Déterminer l'image réciproque par $f$ de $\{1\}$.}
\reponse{
Par définition, $f$ prend deux valeurs : $1$ et $-1$. Donc 
\[f(\Z)=\{-1;1\}\]
Enfin, l'image réciproque de $1$ est composé de tous les éléments de $\Z$ ayant pour image $1$, c'est-à-dire les nombres pairs. Ainsi, $\ds{f^{-1}(\{1\}) = 2\Z}$.
}
}

%%% Fin exercice %%%
